\textbf{Correspondence status: not reviewed. The general kernel, operator, and trace modules compile with the pinned Lean and Mathlib versions. Compilation and the chapter's axiom audit are separate from Steven's review of the mathematical correspondence.}

The formalization uses Mathlib's actual complex $L^2$ equivalence classes, Bochner integrals, continuous linear maps, orthogonal complements, bases, and endomorphism trace. The book's finite family is indexed by $\operatorname{Fin}(N)$, including $N=0$. Several shared results and their section specializations hold for arbitrary measures: sigma-finiteness is not needed for these finite-sum identities. They therefore apply to the manuscript's sigma-finite setting without an additional hypothesis. The ambient $L^2$ space is not assumed finite-dimensional.

For \hyperref[finite-orthonormal-coefficients]{orthonormal coefficients}, the shared proof follows the prose: subtract the finite coefficient sum, show that the remainder is orthogonal to the span, and prove uniqueness of both remainder and coefficients. Applying the given linear operator yields the spectral expansion. Mathlib's inner product is linear in its second argument, so \texttt{bookInner f g} reverses those arguments. The integral theorem proves that it equals the book's $\int_X f\overline g\,d\mu$. The finite matrix operator is a sum of actual rank-one continuous linear maps; its pointwise expansion is proved, rather than included as an assumed property. The operator norm estimate, its image in $V$, its zero action on $V^\perp$, and its action on the basis vectors are all separate consequences.

For the \hyperref[finite-spectral-kernel]{spectral kernel} and \hyperref[finite-kernel-operator]{general matrix kernel}, products of two $L^2$ representatives are proved integrable using Mathlib's $L^2$ theory. The finite integral computation is then related almost everywhere to the actual $L^2$ operator. The spectral specialization identifies the given linear $T$ with the diagonal finite matrix operator, using its eigenvector and orthogonal-complement hypotheses; boundedness of $T$ is a conclusion, not an assumption. The same specialization gives its diagonal-integral and eigenvalue-sum trace identity.

The implementation of the \hyperref[finite-kernel-trace-definition]{finite-rank trace} makes one explicit change of proof route. It defines the trace using Mathlib's trace of the restriction to the actual range of the operator. The theorem \texttt{finiteRankTrace\_eq\_supported\_trace} proves equality with the trace on every finite-dimensional supporting subspace $W$. It factors the operator through its range and uses cyclicity for maps between finite-dimensional spaces, whereas the prose compares $W$ and $W'$ through their common containing space. This proof actually needs only that $W$ contains the image; vanishing on $W^\perp$ remains a hypothesis where the manuscript uses it to reconstruct the spectral operator. Thus the canonical definition has the same value as the manuscript's definition, with the difference in construction and proof made available for Steven's review.

For the \hyperref[finite-kernel-trace-identity]{general trace identity}, the actual restriction to the span is represented in the basis built from the orthonormal family. Its matrix is proved to be $A$. The diagonal integral, the canonical operator trace, and the sum of the characteristic polynomial's root multiset are then equated. Repeated roots give algebraic multiplicity, and no diagonalizability assumption is made. Taking $A=I_N$ gives the actual orthogonal projection onto the finite span and trace $N$. On a finite counting-measure space, a separate point-indicator calculation identifies the kernel with the matrix entries and its integral action with matrix-vector multiplication.

The initial $L^2$ integral formulas use Mathlib's canonical representatives. Explicit almost-everywhere bridges allow them to be replaced by any chosen representatives of the same classes. Both the operator integral and the specified diagonal integral are covered, as is integrability of the replaced diagonal. For \hyperref[kernel-diagonal-representatives]{representative independence}, equality of the finitely many one-variable representatives gives base-measure almost-everywhere equality of the diagonals separately from product-measure equality of the kernels. The counterexample is formalized for actual Lebesgue measure restricted to $[0,1]$: the diagonal indicator equals zero almost everywhere for the product measure, every row integral against every input is zero, but its diagonal integral is $1$, while the zero kernel's is $0$.

The concrete \hyperref[finite-kernel-rank-one-example]{rank-one example} uses actual Lebesgue measure on $[0,1]$ and the $L^2$ class of $u(x)=\sqrt{3}x$. Its norm is proved to be $1$ by evaluating the integral of $3x^2$. The formalization proves the eigenvector and orthogonal-complement statements, the chosen kernel formula $K(x,y)=3\alpha xy$, the almost-everywhere integral formula for the actual operator, and equality of its diagonal integral and trace with $\alpha$. These calculations are exposed independently in the shared rank-one module and specialized in this section's companion.

No admitted proof or project axiom supplies these identities. The pinned dependency and build records, declaration-level axiom output, and exact source versions belong to the chapter audit. Human correspondence review remains pending, and no video-production approval follows from these machine checks.
